\subsection{The boundary norm equations and original factorial exponents}
\label{sec:ghost-boundary-norms}

\begin{proposition}[Hasse norm principle for the three boundary fields]
\label{prop:ghost-boundary-hasse-norm}
Each algebra \(K_i\) of (25) is a field of degree nine.
Its normal closure has group
\(G=H_3\rtimes C_2\), of order fifty-four, with
the involution acting by minus one on the two
abelian Heisenberg generators and trivially on its
centre. A point stabilizer is \(S_3\).
Each \(K_i/\mathbb Q\) satisfies the Hasse norm
principle: a rational number that is a norm at
every completion is an actual global norm.
Its norm-one torus also satisfies weak approximation.
\end{proposition}
\begin{proof}
First consider \(D_0\). Relative to the rational
tangent at zero, the geometric cover is the
characteristic class-two exponent-three quotient
\(H_3\) of the free geometric fundamental group.
The arithmetic action sends its two generators
\(x,y\) to \(x^\chi,y^\chi\). Indeed the
commutator-valued conjugating term in the actual
tangential action becomes central in this quotient.
Its action therefore factors through
\(\chi\bmod3\in C_2\).
The boundary fibre of the pointed twist at three
is the actual coset set \(H_3/\langle x\rangle\),
with the affine Galois action given by the
point-three cocycle and this cyclotomic action.
The tangential reference sheet at zero is Galois
fixed before twisting, which fixes this coset
comparison without a residual boundary translation.

Over \(\mathbb Q(\mu_3)\), the abelian image of
the point-three cocycle is all of \(\mathbb F_3^2\):
the Kummer classes of \(3,-2\) are independent
by their valuations and the injectivity (22).
Any subgroup of \(H_3\) surjecting onto its
abelianization contains the nonzero commutator
of two independent lifts, hence contains the
whole centre and is \(H_3\). Thus the affine
Galois image is all of \(G\). Its stabilizer of
the identity coset is
\(\langle x\rangle\rtimes C_2\simeq S_3\).
That subgroup has trivial core: its intersection
with \(H_3\) is a noncentral order-three subgroup
whose conjugates intersect trivially, and an
element of the core outside \(H_3\) would commute
with \(H_3\), since its commutators lie in the
already trivial intersection of the core with
\(H_3\). This contradicts its minus-one action
on the abelianization. The action is therefore faithful
and transitive of degree nine.

For the other two fibres use the actual automorphisms
\(t\mapsto1-t\) and \(t\mapsto1/t\) of \(U\).
They carry the pointed cover to the same
characteristic quotient cover pointed at \(-2\)
and \(1/3\), respectively, and carry the relevant
boundary to zero. This follows functorially from
the characteristic geometric kernel and the
image of the rational point section; no choice
of a Galois-invariant path is required.
The corresponding abelian pairs are \((-2,3)\)
and \((1/3,2/3)\). Their cube classes are again
independent: their valuation matrices at two
and three have nonzero determinant modulo three.
The preceding coset argument applies verbatim
with the transported generators and gives the
same group and stabilizer for each fibre.

We now compute the norm obstruction. The Heisenberg
group has the presentation
\[
 H_3=\langle x,y\mid x^3=y^3=1,
       [x,[x,y]]=[y,[x,y]]=1\rangle.                 \tag{36}
\]
The last two relations make \([x,y]\) central;
\([x^3,y]=1\) then gives \([x,y]^3=1\).
Collection leaves at most twenty-seven words,
and the standard Heisenberg matrices give
twenty-seven distinct words, proving (36).

Consider a central extension of \(H_3\) by
\(\mathbb F_3\), and lift \(x,y\) to \(X,Y\).
The values of the four relators in (36) lie in
that central kernel. They determine the extension
class: when all four vanish, the lifted generators
define a splitting by (36). Changing the lifts
by central kernel elements leaves the four values
unchanged. Pullback by the involution
\(x\mapsto x^{-1},y\mapsto y^{-1}\) negates
every one of these four values. For the last two,
the extension has class at most three, and the
weight-three commutator changes sign under this
substitution; all weight-four errors vanish.
Thus \(C_2\) acts by minus one on
\(H^2(H_3,\mathbb F_3)\).
Restriction from \(G\) to \(H_3\) is injective
on \(\mathbb F_3\)-cohomology, since its composite
with corestriction is multiplication by two;
its image is invariant. Consequently
\[
 H^2(G,\mathbb F_3)=0,
             \qquad H^3(G,\mathbb Z)\{3\}=0.        \tag{37}
\]
For the second conclusion, use
\(H^3(G,\mathbb Z)=H^2(G,\mathbb Q/\mathbb Z)\)
and the multiplication-by-three exact sequence;
a nonzero finite three-primary group would have
nonzero three-torsion detected from
\(H^2(G,\mathbb F_3)\).
Restriction to the order-two complement is injective
on the two-primary part of \(H^3(G,\mathbb Z)\),
because its index is twenty-seven. The target
\(H^3(C_2,\mathbb Z)\) is zero by the cyclic
resolution. These are the only possible primary
parts for a group of order fifty-four, so in fact
\(H^3(G,\mathbb Z)=0\).

For \(H\simeq S_3\) the character group of the
norm-one torus of \(K_i/\mathbb Q\) is
\[
 X_i=\operatorname{coker}
       (\mathbb Z\longrightarrow\mathbb Z[G/H]),
 \qquad 1\longmapsto\sum_{G/H}gH.                   \tag{38}
\]
The long cohomology sequence and Shapiro give
\[
 H^2(G,\mathbb Z)\longrightarrow
 H^2(H,\mathbb Z)\longrightarrow H^2(G,X_i)
                  \longrightarrow H^3(G,\mathbb Z).
\]
Here \(H^2(H,\mathbb Z)=\operatorname{Hom}
(H,\mathbb Q/\mathbb Z)\simeq\mathbb Z/2\).
Its generator is the restriction of the character
\(G\to C_2\subset\mathbb Q/\mathbb Z\), so the
first displayed map is surjective. Consequently
\(H^2(G,X_i)=0\).
Tate's norm-torus duality, in the precise form
\cite[Theorem 2.3, equation (2.6), and Definition 2.5]
{GhostMacedoNewton2022}, embeds the dual of the
norm Hasse obstruction into \(H^2(G,X_i)\).
It is therefore zero. Moreover, Theorem 2.2 and
equation (2.4) of the same source identify the
Picard cohomology of a smooth compactification of
the norm-one torus with a subgroup of this same
zero group. Voskresenski\u\i's exact sequence,
Theorem 2.1 there, then makes its weak-approximation
defect zero. This proves both assertions. The
independent normalized-bar-complex calculation in
\path{code/check_boundary_norm_group.py}
checks the invariant cohomology in (37), as well as
the group order, stabilizers, and faithful coset actions.
\end{proof}

\begin{proposition}[An unscaled A.44 family surviving the affine-cover abelian quotient]
\label{prop:ghost-original-genjac-survival}
The original recursion with exponents \((j+4)!\)
admits ordinary infinite-support outputs with
\[
 q_j\equiv r_j\equiv1\pmod {27},\qquad
 3\notin(\mathbb F_{q_j}^\times)^3,
       \quad2\notin(\mathbb F_{r_j}^\times)^3,        \tag{39}
\]
and an actual global continuous section of
\(\pi_1(U)/\overline{[K,K]}\) lying above that
full completed pair. It factors through the
generalized-Jacobian quotient of the pointed
cover \(Y\), and its compact Tate-module class
is zero at every coefficient prime.
\end{proposition}
\begin{proof}
At the place three in each A.44 induction stage,
choose \(\xi_3=1\). The old equation then has
a near-one solution for \(\eta_3\): its derivative
has valuation \(e=v_3(n_j)\), while its error has
valuation at least
\(H_j(3)=2v_3(n_{j+1})+j+3\).
The same Hensel calculation as in A.44 gives
\[
 v_3(\eta_3-1)\ge H_j(3)-e\ge4,
       \qquad v_3(\eta_3-1)\ge v_3(n_{j+1}).         \tag{40}
\]
The first root is exactly one at the initial
stage. The existing CRT modulus has three-adic
valuation at least four. Thus both new primes
are one modulo twenty-seven, while all original
congruences, including A.44.14, are retained.
The field-intersection argument of
Proposition~\ref{prop:ghost-elliptic-passing}
adds the two noncube conditions without changing
these prescribed CRT automorphisms. Induction
constructs (39) within the original A.44 recursion.

We prove that each \(q_j\) is an actual global
norm from \(K_0\), and each \(r_j\) from \(K_1\).
First, these fields are unramified outside two
and three. Indeed (18) identifies
\[
 \begin{gathered}
 D_0=[3]^{-1}(-P_0-T),\qquad
 D_1=[3]^{-1}(-P_0+T),\\
 D_\infty=[3]^{-1}(-P_0).
 \end{gathered}                                      \tag{41}
\]
At a prime greater than three, \(E'\) has good
reduction and multiplication by three is finite
etale on its proper smooth model. The inverse
image of each integral rational section in (41)
is finite etale over the valuation ring, proving
the asserted unramifiedness.

At \(q=q_j\), the proof of
Proposition~\ref{prop:ghost-elliptic-passing} gives
\(E'(\mathbb F_q)\{3\}=\langle T\rangle\) and
\(P_0+T\) of prime-to-three order. Therefore
\(-P_0-T\) is three-divisible in
\(E'(\mathbb F_q)\), and etale Hensel lifting in
(41) supplies a \(\mathbb Q_q\)-point of \(D_0\).
Thus \(K_0\otimes\mathbb Q_q\) has a factor
\(\mathbb Q_q\); its product norm is surjective.
The analogous argument gives a degree-one factor
of \(K_1\otimes\mathbb Q_r\) at \(r=r_j\).

At every other prime greater than three the
specified rational prime \(q_j\), or \(r_j\),
is a unit, and units are norms from each
unramified local extension. At two an odd unit
is a ninth power in \(\mathbb Q_2\); at three
every unit in \(1+27\mathbb Z_3\) is a ninth
power, by the logarithm isomorphism on
\(1+3\mathbb Z_3\). Such a ninth power is a
norm through the diagonal embedding into a
degree-nine etale algebra. At the real place
there is a real factor since the degree is odd,
so the product norm is surjective there as well.
These checks give actual local norms at all
places, not merely norms in finite coefficients.
Proposition~\ref{prop:ghost-boundary-hasse-norm}
therefore supplies
\[
 z_j\in K_0^\times,\quad w_j\in K_1^\times,
              \qquad Nz_j=q_j,\quad Nw_j=r_j.       \tag{42}
\]

The products
\[
 \alpha_0=\prod_j z_j^{n_j/3},\qquad
 \alpha_1=\prod_j w_j^{n_j/3},\qquad
                    \alpha_\infty=1                 \tag{43}
\]
converge in the three-adic multiplicative completions:
\(v_3(n_j/3)\) tends to infinity, so every tail
is zero at each fixed coefficient precision.
Their norms are exactly
\(\prod q_j^{n_j/3}\) and \(\prod r_j^{n_j/3}\).
Formula (27) with \(a=0\) produces a global
three-adic section of the generalized-Jacobian
quotient above the asserted pair, with proper
class zero.

For every coefficient prime \(\ell\ne3\), the
block-constant map (32) gives the same conclusion
for an arbitrary pair, since multiplication by
twenty-seven is an automorphism of
\(H^1(G_{\mathbb Q},\mathbb Z_\ell(1))\).
Choose these explicit toric cocycles, and take
their product with the three-adic cocycle of
(43). A map to the product Tate module is
continuous exactly when each coordinate is
continuous. It is therefore an actual global
cocycle in \(T\mathcal J\), and the rational
origin section gives the claimed full profinite
section. Its proper class is zero in every
coordinate. The map into the downstairs quotient
is precisely the inclusion in
Proposition~\ref{prop:ghost-genjac-image}.
\end{proof}

This supplies a single original-factorial ghost
surviving the displayed infinite affine-cover
abelian quotient. The kernel still retained in
the full cover is \(\overline{[K,K]}\). Neither
its lifting problem nor its local ordinary-point
comparison at that nonabelian level is settled by
(42)--(43); those are the remaining constraints
on this explicitly constructed family.
